% 第6章第9节：本章小结。

\section{本章小结}
\label{sec:ls-analysis-summary}
\label{sec:linear-algebra-summary}
\label{sec:matrix-analysis-summary}

本章先研究线性作用在什么样的基下最简单，再用这些简单表示解决长度约束下的极值、低秩逼近和伸缩分离三个问题，最后研究误差、扰动与反复作用。

前五节逐步调整对基的要求。由特征向量组成的基使矩阵成为对角形。特征值是特征多项式$p_{\bm A}(z)=\det(z\bm I-\bm A)$的根，不随基改变；Cayley--Hamilton定理说明$p_{\bm A}(\bm A)=\bm0$。$n$阶矩阵可对角化，当且仅当它有$n$个线性无关的实特征向量，这时$\bm A^k=\bm V\bm\Lambda^k\bm V^{-1}$。要求基标准正交时，能对角化的恰好是实对称矩阵，分解可写成$\bm S=\sum_\lambda\lambda\bm P_\lambda$，其中谱投影由$\bm S$唯一确定。一般方阵只能化为实Schur形，正规矩阵则能化为分块对角形。输入、输出各取一组标准正交基时，任何矩阵都有奇异值分解，由它读出谱范数、Frobenius范数和秩。

随后三节各解决一个具体问题。在单位长度下，二次型$\bm x^\top\bm S\bm x$的最大、最小值是$\bm S$两端的特征值；把长度换成$\bm x^\top\bm B\bm x=1$，用$\bm B=\bm L^\top\bm L$换元就化为普通情形，极值由广义特征方程$\bm S\bm v=\mu\bm B\bm v$给出。秩不超过$k$的近似矩阵中，截断奇异值分解$\bm A_k$使重建误差最小：先把选矩阵化为选子空间，秩一时用Rayleigh商求出最优直线，秩$k$时用迹极值证明最优性，两种范数下的最小误差分别是$(\sum_{i>k}\sigma_i^2)^{1/2}$与$\sigma_{k+1}$。把方阵写成$\bm A=\bm W\bm H$时，先由$\bm W^\top\bm W=\bm I$消去正交因子，得到$\bm H^2=\bm A^\top\bm A$，沿谱开方求出唯一的半正定$\bm H$，可逆时再由$\bm W=\bm A\bm H^{-1}$求出正交因子；这个$\bm W$也是Frobenius范数下离$\bm A$最近的正交矩阵。开方与求逆都是把每个特征值换成函数值，由此得到沿谱定义的矩阵函数$f(\bm S)$。

最后三节读出三类彼此相关、但不能互相代替的性质。第一类是一次作用对输入误差的放大：绝对误差的最大放大倍数是最大奇异值，固定输入处的相对放大由相对条件数刻画，对全部输入取最坏值得到$\kappa_2=\sigma_{\max}/\sigma_{\min}$。第二类是矩阵本身受扰动时的变化：对称矩阵的有序特征值与任意矩阵的奇异值移动不超过$\lVert\bm E\rVert_2$，可据此保证正定性和秩；可逆矩阵到奇异矩阵的距离是$\sigma_{\min}$；特征子空间的转角由扰动与谱隙之比控制；一般方阵的特征值还受特征基条件数影响，缺少特征向量时可能很敏感。第三类是固定更新$\bm x_{k+1}=\bm T\bm x_k+\bm c$的长期行为：初值影响消失当且仅当$\rho(\bm T)<1$，同一条件使Neumann级数收敛，所有初值趋于同一平衡$(\bm I-\bm T)^{-1}\bm c$；谱半径决定长期速度，非正规矩阵可能先放大、再衰减，每步的有界误差累积成有界的误差水平。

第~\ref{chap:linear-equations-and-optimization}章求解线性方程、最小二乘与矩阵方程时，将用到奇异值分解、Rayleigh商、条件数、谱、Cayley--Hamilton定理、实Schur形与Neumann级数。本章的误差结论都是有限变化界；第~\ref{chap:multivariate-calculus}章再对一般矩阵映射建立微分，求局部条件数与一阶响应。
